\documentclass[11pt]{article}
\usepackage[a4paper,margin=25mm]{geometry}
\usepackage{amsmath,amssymb,amsthm,mathtools}
\usepackage[hidelinks]{hyperref}
\allowdisplaybreaks[2]
\setlength{\emergencystretch}{2em}
\title{A Nordhaus--Gaddum bound for normalized algebraic connectivity}
\author{Leo Gao\\[-1pt]{\small AI-assisted research}}
\date{November 26, 2026 \quad$\cdot$\quad Research draft}
\begin{document}
\maketitle
\begin{abstract}
We derive an absolute $c>0$ such that the sum of the second normalized Laplacian eigenvalues of a graph and its complement is at most $2-c/sqrt n$. A family with deficit $(1+o(1))/sqrt n$ gives the sharp order. For regular graphs, we obtain the sharp bound $2-2/(sqrt n+1)$, attained by Paley graphs.
\end{abstract}
\section{The conjecture}

On June 10, 2026, Mark Kempton, Xavier Zaitzeff and Sibi Muthuprakash released their research paper:

\emph{Nordhaus--Gaddum upper bounds for graph connectivity parameters}~\cite{KZM}.

In Section 4.2, Conjecture 2 proposes the asymptotic inequality

\[
\lambda_2(G)+\lambda_2(\overline G)
\leq 2-O(n^{-1/2}).
\]

Here, \(\lambda_2\) is the second-smallest eigenvalue of the normalized graph Laplacian, and \(\overline G\) is the graph complement.

I interpret the conjecture in its intended, nontrivial sense: there exists a strictly positive universal constant \(c\) such that the deficit from 2 is at least \(c/\sqrt n\).

\section{Main theorem}

\textbf{Theorem.} There exists an absolute constant \(c>0\) such that, for every finite simple graph \(G\) on \(n\geq 3\) vertices,

\[
\boxed{
\lambda_2(G)+\lambda_2(\overline G)
\leq 2-\frac{c}{\sqrt n}.
}
\]

Furthermore, the order \(n^{-1/2}\) is optimal: there exists a sequence of graphs satisfying

\[
\boxed{
\lambda_2(G_n)+\lambda_2(\overline G_n)
=2-\frac{1+o(1)}{\sqrt n}.
}
\]

The first statement resolves the conjecture. The second establishes its asymptotic sharpness and improves the leading coefficient achieved by the constructions motivating the original conjecture.

\section{Proof}

\subsection{Notation and initial reduction}

Let \(A\) be the adjacency matrix of \(G\), \(D\) its diagonal degree matrix, \(d\) its degree vector, and

\[
s=\sum_{i=1}^n d_i.
\]

Use barred notation for corresponding quantities belonging to \(\overline G\).

We have the identities

\[
A+\overline A=J-I,
\]

\[
d+\overline d=(n-1)\mathbf1.
\]

For connected graphs, define the normalized adjacency matrix

\[
S=D^{-1/2}AD^{-1/2}.
\]

Its largest eigenvalue is 1, and its second-largest eigenvalue is

\[
a=1-\lambda_2(G).
\]

Similarly, let

\[
b=1-\lambda_2(\overline G).
\]

We use the normalized Laplacian convention in which an isolated vertex contributes a zero eigenvalue. If either graph is disconnected, its second normalized Laplacian eigenvalue is zero. The other graph is connected and has normalized Laplacian trace $n$, so its second eigenvalue is at most $n/(n-1)$. This gives the result in this case.

We may therefore assume that both graphs are connected. Neither can be complete, so interlacing on two nonadjacent vertices gives

\[
a,b\geq0.
\]

It suffices to prove that

\[
a+b=\Omega(n^{-1/2}).
\]

Suppose, for a contradiction, that there exists a sequence of graphs with \(n\to\infty\) such that

\[
\varepsilon:=a+b=o(n^{-1/2}).
\tag{1}
\]

In particular, both normalized algebraic connectivities tend to 1.

\subsection{The central matrix inequality}

Define two symmetric matrices:

\[
P=aD+(1-a)\frac{dd^\top}{s}-A,
\]

\[
Q=b\overline D+
(1-b)\frac{\overline d\,\overline d^\top}{\overline s}
-\overline A.
\]

\textbf{Lemma 1.} Both \(P\) and \(Q\) are positive semidefinite.

\textbf{Proof.} The unit eigenvector of \(S\) belonging to its largest eigenvalue is

\[
u=\frac{D^{1/2}\mathbf1}{\sqrt s}.
\]

All other eigenvalues of \(S\) are at most \(a\). Consequently,

\[
aI+(1-a)uu^\top-S\succeq0.
\]

Congruence by \(D^{1/2}\) gives \(P\succeq0\). The complementary argument gives \(Q\succeq0\). \(\square\)

Choose a unit vector \(q\) perpendicular to \(\mathbf1\) such that

\[
d\in\operatorname{span}\{\mathbf1,q\}.
\]

If \(G\) is regular, any such unit vector suffices. Because \(d+\overline d=(n-1)\mathbf1\), the complementary degree vector also belongs to this span.

Let \(\Pi\) denote orthogonal projection onto

\[
U=\operatorname{span}\{\mathbf1,q\}^{\perp}.
\]

For every \(x\in U\), the degree-vector terms vanish, giving

\[
x^\top(P+Q)x
=x^\top(I+aD+b\overline D)x.
\]

Therefore, with

\[
K=1+(n-1)(a+b),
\]

positive semidefiniteness implies

\[
0\leq x^\top Px\leq K\|x\|^2.
\tag{2}
\]

Write \(r_i\) for the adjacency row corresponding to vertex \(i\).

By the Cauchy--Schwarz inequality for positive semidefinite matrices,

\[
|e_i^\top Px|^2
\leq P_{ii}\,x^\top Px.
\]

But

\[
P_{ii}
=a d_i+(1-a)\frac{d_i^2}{s}
\leq a d_i+\frac{d_i^2}{s}.
\]

For \(x\in U\),

\[
e_i^\top Px=a d_i x_i-r_i^\top x.
\]

Combining these expressions and maximizing over unit vectors in \(U\) yields

\[
\|\Pi r_i\|
\leq a d_i+
\sqrt{K\left(a d_i+\frac{d_i^2}{s}\right)}.
\]

Squaring and dividing by \(d_i>0\), we obtain the central inequality:

\[
\boxed{
\frac{\|\Pi r_i\|^2}{d_i}
\leq
2a^2d_i+2K\left(a+\frac{d_i}{s}\right).
}
\tag{3}
\]

The identical inequality holds for complementary adjacency rows with \(a,d,s\) replaced by \(b,\overline d,\overline s\).

\subsection{The structural consequence}

Averaging (3) over all vertices gives

\[
\frac1n\sum_i
\frac{\|\Pi r_i\|^2}{d_i}
\leq
2a^2\frac{s}{n}
+2K\left(a+\frac1n\right).
\]

Under assumption (1), the right-hand side tends to zero. Thus

\[
\boxed{
\frac1n\sum_i
\frac{\|\Pi r_i\|^2}{d_i}=o(1).
}
\tag{4}
\]

Moreover, whenever \(s=\Theta(n^2)\), inequality (3) gives the stronger, uniform conclusion

\[
\boxed{
\max_i\frac{\|\Pi r_i\|^2}{d_i}=o(1).
}
\tag{5}
\]

Both conclusions apply to the complement as well.

This is the main mechanism of the proof: if the conjecture were false, the normalized adjacency rows of both graphs would become almost entirely concentrated in the same two-dimensional space.

The binary nature of adjacency matrices makes that condition extraordinarily restrictive.

We now derive a contradiction in every possible asymptotic density regime.

\subsection{Case I: Both graphs have nonvanishing edge density}

Suppose, along a subsequence,

\[
\frac{s}{n^2}\longrightarrow p,
\qquad 0<p<1.
\]

Then both graphs satisfy the uniform estimate (5).

There must exist vertices with degree bounded away from both 0 and \(n\), proportionally to \(n\).

Otherwise, along a subsequence, all vertices would have degree either \(o(n)\) or \(n-o(n)\). Since \(p\in(0,1)\), both groups would contain a positive proportion of all vertices. Edges between these groups would simultaneously have to be almost absent, by the low degrees, and almost complete, by the high degrees. This is impossible.

Choose such a vertex and let \(H\) be its neighborhood, with

\[
|H|=\Theta(n),
\qquad |H^c|=\Theta(n).
\]

Define the centered unit indicator

\[
w_H=
\frac{\mathbf1_H-(|H|/n)\mathbf1}
{\sqrt{|H|(1-|H|/n)}}.
\]

By (5), this vector lies within \(o(1)\) of the one-dimensional space spanned by \(q\). Reversing the sign of \(q\) if necessary,

\[
\|q-w_H\|=o(1).
\tag{6}
\]

Every coordinate of \(w_H\) is \(O(n^{-1/2})\).

Now suppose a vertex had degree \(d_i=o(n)\). By (5), the unit vector

\[
\frac{r_i}{\sqrt{d_i}}
\]

would have negligible projection onto \(U\). Its projection onto the constant-vector direction also has norm

\[
\sqrt{d_i/n}=o(1).
\]

It must therefore have inner product of absolute value \(1-o(1)\) with \(q\).

But (6) implies

\[
\left|
\left\langle
\frac{r_i}{\sqrt{d_i}},q
\right\rangle
\right|
\leq O\left(\sqrt{\frac{d_i}{n}}\right)+o(1)
=o(1),
\]

a contradiction.

Applying the same argument to \(\overline G\) excludes vertices of complementary degree \(o(n)\).

Consequently, all degrees in both graphs are bounded below by a fixed positive proportion of \(n\), for sufficiently large \(n\) along this sequence.

By (5) and (6), every adjacency row is within Euclidean distance \(o(\sqrt n)\), uniformly, of a vector constant on \(H\) and on \(H^c\).

Rounding these two constants to their nearest binary values changes at most \(o(n)\) entries: each mismatching binary entry contributes at least \(1/4\) to the squared approximation error.

The all-zero and all-one patterns are excluded by the degree bounds. Thus every adjacency row differs in only \(o(n)\) entries from either

\[
\mathbf1_H
\quad\text{or}\quad
\mathbf1_{H^c}.
\]

Write \(h_i=1\) on \(H\) and \(h_i=-1\) outside \(H\). For each row, let \(t_i\in\{-1,1\}\) identify which pattern approximates it.

Then \(A\) differs in \(o(n^2)\) entries from the matrix

\[
M_{ij}=\frac{1+t_i h_j}{2}.
\]

Since \(A\) is symmetric, \(M\) differs from its transpose in only \(o(n^2)\) entries.

However,

\[
M_{ij}\ne M_{ji}
\iff
t_i h_i\ne t_j h_j.
\]

If \(k\) is the number of vertices for which \(t_i h_i=1\), the number of ordered asymmetric entries is exactly

\[
2k(n-k).
\]

Hence

\[
k(n-k)=o(n^2).
\]

It follows that all but \(o(n)\) vertices have the same value of \(t_i h_i\).

Therefore \(G\) is within \(o(n^2)\) adjacency changes of either:

\begin{itemize}
\item two disjoint cliques on \(H\) and \(H^c\); or
\item a complete bipartite graph between \(H\) and \(H^c\).

\end{itemize}

In either case, one of \(G,\overline G\), denoted by \(F\), is almost a union of two macroscopic cliques.

For any nonempty proper vertex set \(T\), the normalized Laplacian Rayleigh quotient gives

\[
\lambda_2(F)\leq
|\partial_F T|
\left(
\frac1{\operatorname{vol}_F(T)}
+
\frac1{\operatorname{vol}_F(T^c)}
\right).
\tag{7}
\]

Here, \(\operatorname{vol}_F(T)\) is the sum of degrees in \(T\), and \(\partial_F T\) is its edge boundary.

Taking \(T=H\), we have

\[
|\partial_F H|=o(n^2),
\]

while both volumes are \(\Theta(n^2)\).

Equation (7) gives

\[
\lambda_2(F)=o(1).
\]

This contradicts the assumption that both normalized algebraic connectivities converge to 1.

\subsection{Case II: One graph is asymptotically sparse}

It remains to handle the regime in which one graph has vanishing edge density. Exchanging the roles of the graphs if necessary, assume

\[
s=o(n^2).
\]

The complement is consequently dense.

By (4) and the vanishing average degree fraction, there exists a set \(W\) containing \(n-o(n)\) vertices such that, uniformly for \(i\in W\),

\[
d_i=o(n),
\qquad
\frac{\|\Pi r_i\|^2}{d_i}=o(1).
\]

For each such vertex, its normalized adjacency row has negligible projections onto both \(U\) and the constant-vector direction. Therefore

\[
\frac{r_i}{\sqrt{d_i}}
=\pm q+o(1)
\]

in Euclidean norm.

All these normalized rows are nonnegative vectors. Two such vectors cannot approach opposite unit vectors, because their inner product is nonnegative. Thus their signs agree, and we may orient \(q\) so that

\[
\frac{r_i}{\sqrt{d_i}}
=q+o(1)
\tag{8}
\]

uniformly over \(W\).

Choose one vertex in \(W\), let \(H\) be its neighborhood, and put \(h=|H|\). Then

\[
1\leq h=o(n)
\]

and

\[
q=\frac{\mathbf1_H}{\sqrt h}+o(1).
\tag{9}
\]

Combining (8) and (9),

\[
\frac{|N_G(i)\cap H|}{\sqrt{d_i h}}
=1-o(1).
\]

Since the intersection is at most \(\min(d_i,h)\), this forces, uniformly over \(W\),

\[
d_i=(1+o(1))h,
\]

\[
|N_G(i)\cap H|=(1-o(1))h.
\]

Almost every vertex is therefore adjacent to almost every vertex of \(H\).

Summing these incidences yields

\[
\sum_{j\in H}d_j\geq(1-o(1))nh.
\]

Consequently,

\[
\sum_{j\in H}\overline d_j=o(nh).
\tag{10}
\]

By (10), there exists a subset \(H'\subseteq H\) with

\[
|H'|=(1-o(1))h
\]

such that

\[
\overline d_j=o(n)
\]

uniformly for \(j\in H'\).

Now apply the dense-graph estimate (5) to the complement. Each complementary normalized adjacency row belonging to \(H'\) must approach either \(q\) or \(-q\).

Again, nonnegativity and (9) exclude the negative sign. Therefore

\[
\frac{\overline r_j}{\sqrt{\overline d_j}}
=
\frac{\mathbf1_H}{\sqrt h}+o(1).
\]

As before, the binary intersection argument gives

\[
\overline d_j=(1+o(1))h
\]

and

\[
|N_{\overline G}(j)\setminus H|=o(h),
\]

uniformly for \(j\in H'\).

It follows that

\[
\operatorname{vol}_{\overline G}(H')
=(1+o(1))h^2.
\]

The number of complementary edges leaving \(H'\) satisfies

\[
|\partial_{\overline G}H'|
\leq
|H'|\,o(h)+|H'||H\setminus H'|
=o(h^2).
\]

Since the complement is dense and \(h=o(n)\), the remaining vertices have volume \(\Theta(n^2)\).

Using (7) again, we conclude that

\[
\lambda_2(\overline G)=o(1).
\]

This contradicts \(\lambda_2(\overline G)\to1\).

\subsection{Conclusion of the proof}

Every possible limiting edge-density regime produces a contradiction.

There can therefore be no sequence of graphs satisfying

\[
\sqrt n\left[
2-\lambda_2(G)-\lambda_2(\overline G)
\right]\longrightarrow0.
\]

Equivalently,

\[
\liminf_{n\to\infty}
\min_{|V(G)|=n}
\sqrt n\left[
2-\lambda_2(G)-\lambda_2(\overline G)
\right]>0.
\]

This establishes the existence of a universal \(c>0\) for all sufficiently large \(n\).

For the remaining finite orders \(n\geq3\), the deficit is strictly positive. To see this in the only nontrivial case, suppose both graphs are connected and \(a=0\). For nonadjacent vertices \(i,j\), set

\[
x=\frac{e_i}{d_i}-\frac{e_j}{d_j}.
\]

Then \(x^\top Px=0\), so positive semidefiniteness forces \(Px=0\), implying that \(i,j\) have identical neighborhoods. Nonadjacency thus partitions \(G\) into independent classes, with every possible edge between distinct classes. Such a connected complete multipartite graph has a disconnected complement, a contradiction.

Reducing \(c\) to accommodate finitely many smaller graphs completes the proof.

\[
\boxed{\text{Conjecture established.}\quad\blacksquare}
\]

\section{Asymptotic sharpness: an improved extremal construction}

Proving the bound is only half of the extremal problem. Can the exponent or leading coefficient be improved?

Consider the graph \(G_{k,m}\) constructed from a clique on \(k\) vertices by attaching \(m\) pendant vertices to distinct clique vertices, where \(2\leq m\leq k\).

Set

\[
n=k+m,\qquad m\to\infty,\qquad m=o(n).
\]

Both \(G_{k,m}\) and its complement are connected.

In \(G_{k,m}\), attached clique vertices have degree \(k\), and the pendant vertices have degree 1.

The normalized adjacency matrix decomposes into a positive rank-one clique term, a negative diagonal term, and a weighted matching term whose operator norm is \(1/\sqrt{k}\).

By the eigenvalue minimax principle,

\[
a\leq\frac1{\sqrt k}.
\]

On the contrast subspace between matched vertex pairs, the normalized adjacency matrix acts through the block

\[
\begin{pmatrix}
-1/k & 1/\sqrt k\\
1/\sqrt k & 0
\end{pmatrix}.
\]

Its positive eigenvalue gives the lower bound

\[
a\geq\frac{\sqrt{1+4k}-1}{2k}.
\]

Thus

\[
a=\frac{1+o(1)}{\sqrt n}.
\]

In the complement, the former pendant vertices form a clique and have degree \(n-2\). The former clique vertices form an independent set, with degrees \(m-1\) or \(m\).

Ignoring the deleted matching edges and the negative diagonal correction, the normalized adjacency matrix has at most one positive eigenvalue: its nonzero action is a two-dimensional block with negative determinant.

Restoring the deleted matching perturbs the matrix by an operator of norm

\[
\frac1{\sqrt{(m-1)(n-2)}}.
\]

The negative diagonal correction cannot increase the second-largest eigenvalue. Hence

\[
0\leq b\leq
\frac1{\sqrt{(m-1)(n-2)}}
=o(n^{-1/2}).
\]

Combining these bounds,

\[
\boxed{
a+b=\frac{1+o(1)}{\sqrt n}.
}
\]

This proves the second part of the theorem.

In particular, the optimal asymptotic coefficient is at most $1$.

The original paper's balanced construction has leading coefficient \(\sqrt2\), while its Paley graph construction has leading coefficient 2.

\section{Additional theorem: the sharp regular-graph bound}

The regular-graph case admits a substantially stronger, explicit inequality.

\textbf{Theorem.} For every regular graph \(G\) on \(n\geq3\) vertices with degree between 1 and \(n-2\),

\[
\boxed{
\lambda_2(G)+\lambda_2(\overline G)
\leq 2-\frac{2}{\sqrt n+1}.
}
\]

Equality is attained by Paley graphs.

\textbf{Proof.} Let \(d\) be the regular degree and put \(e=n-1-d\). Denote the adjacency eigenvalues on \(\mathbf1^\perp\) by \(\theta_i\).

Define centered eigenvalues

\[
\xi_i=\theta_i+\frac{d}{n-1}.
\]

The standard adjacency trace identities give

\[
\frac1{n-1}\sum_i\xi_i=0,
\]

\[
\frac1{n-1}\sum_i\xi_i^2
=\frac{nde}{(n-1)^2}.
\]

From the normalized adjacency eigenvalue definitions,

\[
-e\left(b+\frac1{n-1}\right)
\leq\xi_i\leq
d\left(a+\frac1{n-1}\right).
\]

For zero-mean real numbers contained in an interval \([-L,R]\), their mean square is at most \(LR\). This follows by averaging

\[
(x+L)(R-x)\geq0.
\]

Applying this inequality and cancelling \(de\) yields

\[
\left(a+\frac1{n-1}\right)
\left(b+\frac1{n-1}\right)
\geq\frac{n}{(n-1)^2}.
\]

The arithmetic--geometric mean inequality now gives

\[
a+b\geq
\frac{2(\sqrt n-1)}{n-1}
=\frac2{\sqrt n+1}.
\]

Paley graphs attain equality through their two distinct nonprincipal adjacency eigenvalues. \(\square\)

The combination of the regular-graph theorem and the irregular construction demonstrates that \textbf{degree irregularity permits strictly better asymptotic simultaneous connectivity than regularity does}.

\section{Reproduction}
The accompanying \texttt{verify.py} numerically checks the semidefinite and row inequalities on 739 graph-atlas entries and evaluates the scaled deficit of the sharpness family using small spectral blocks.

\section*{Disclosure}
This manuscript was prepared with AI assistance. The arguments and priority claims are open to independent review.

\begin{thebibliography}{9}
\bibitem{KZM} M. Kempton, X. Zaitzeff and S. Muthuprakash, Nordhaus--Gaddum upper bounds for graph connectivity parameters, arXiv:2606.12751 (2026). \href{https://arxiv.org/abs/2606.12751}{Source.}
\end{thebibliography}
\end{document}
